\chapter{Background and Related Work}
\label{ch:background}

This chapter reviews three lines of literature on which the present
work rests: the physics of magnetic winding flux as a flare
precursor (\S\ref{sec:bg-winding}), the self-supervised learning
regime for spatiotemporal data (\S\ref{sec:bg-ssl}), and the
joint-embedding predictive doctrine
(\S\ref{sec:bg-jepa}). It closes with a brief survey of solar
foundation models (\S\ref{sec:bg-surya}) and an explicit
justification of the design choice between joint-embedding and
adversarial inpainting (\S\ref{sec:bg-vs-pathak}).

\section{Magnetic winding flux as a flare precursor}
\label{sec:bg-winding}

The pre-flare configuration of an active region is most directly
observed through its photospheric magnetic field. Two scalar
summaries of that field have received concentrated attention in the
forecasting literature: the unsigned magnetic flux $\Phi = \int
|B_{r}| \, \mathrm{d}A$, and the family of helicity-like topological
invariants that quantify the chirality and entanglement of the field.

Magnetic winding flux, introduced as a forecasting target by Prior
and MacTaggart~\cite{prior2020helicity}, is the second of these.
Unlike the unsigned flux $\Phi$ or the magnetic-helicity flux, the
winding flux is mathematically independent of field strength: it
measures the rate at which topologically twisted flux tubes thread
the photosphere, with a sign that records chirality. Two active
regions with identical~$\Phi$ but opposite handedness yield winding
fluxes of opposite sign.

The forecast relevance of winding flux is established by
Raphaldini~et~al.~\cite{raphaldini2022winding}, who report that
sharp positive inputs of winding flux precede individual flaring
events in AR~11318 by six to seven hours.
Kutsenko~et~al.~\cite{kutsenko2021flare} extend the picture by
linking emergence rate, rather than absolute magnitude, to flare
productivity. Taken together, the literature supports a precursor
mechanism that magnetograms alone do not resolve: rapid input of
\emph{current-carrying} winding marks the photospheric emergence of
topologically twisted flux, and this emergence event precedes the
flare by hours.

The data used in the present work expose the winding flux as a
per-pixel signed pseudoscalar field $W(x, y, t)$ on the
photospheric cutout~$S$, sampled at twelve-minute cadence. The
spatial-mean curve is defined
\begin{equation}
\label{eq:spatial-mean-winding}
\langle W \rangle (t) \;=\;
\frac{1}{|S_{\text{valid}}|}
\int_{S_{\text{valid}}} W(x, y, t) \, \mathrm{d}A,
\end{equation}
restricted to the valid-pixel subset
$S_{\text{valid}} = \{(x, y) : |W(x, y, t)| \le 10^{8}\}$ defined by
the audit of Section~\ref{sec:data-clip}. The bound preserves the
legitimate per-pixel range $10^{5}$--$10^{7}$ established by the
upstream pipeline and excludes the sparse-random numerical glitches
of the source data without discarding real extrema. The
spatial-mean winding-flux curve $\langle W \rangle(t)$ is the
minimum sufficient statistic for the precursor signal observed by
the works cited above; the Stage-1 probe of
Chapter~\ref{ch:probe} targets this curve directly.

\paragraph{Persistence baseline.} The trivial forecaster used as the
lower-bar comparison in Chapter~\ref{ch:probe} is the lag-one
predictor $\hat{y}_{t} = y_{t-1}$. For the wind-flux regression at
twelve-minute cadence, the lag-one predictor exploits the strong
local autocorrelation of $\langle W \rangle(t)$ and is the trivial
floor that any non-trivial probe must beat on a per-cube basis. For
binary flare classification under a $k$-hour forward window, the
analogous baseline is the label-persistence predictor
$\hat{y}_{t} = y_{t-1}$ on the binarised event sequence; under
dense labels this baseline is hard to beat by construction and is
the principal source of optimistic-looking aggregate metrics in
the cross-cube setting.

\section{Self-supervised learning for spatiotemporal data}
\label{sec:bg-ssl}

The dominant self-supervised learning paradigm for image and video
data in the present decade is masked-token pretraining. Three
landmarks define the regime.

The masked autoencoder of He~et~al.~\cite{he2022mae} introduces a
high-mask-ratio (typically 75\%) reconstruction objective in pixel
space for static images, with a small decoder and a large encoder.
The asymmetric encoder--decoder split is the load-bearing design
decision: the encoder processes only visible tokens, and the
decoder is a thin reconstruction head discarded after pretraining.

VideoMAE~\cite{tong2022videomae} extends the masked-autoencoder
paradigm to video by introducing tube-shaped masks that occlude a
patch across all frames simultaneously. The tube mask defeats the
temporal-redundancy shortcut and forces the encoder to learn motion
representations rather than per-frame identity.
VideoMAEv2~\cite{wang2023videomaev2} and
MaskViT~\cite{gupta2022maskvit} subsequently scale the family.
W-MAE~\cite{wmae2023} adapts the regime to atmospheric data.

ClimaX~\cite{nguyen2023climax}, Aurora~\cite{bodnar2024aurora},
Prithvi-WxC~\cite{schmude2024prithviwxc}, and the
SDO-FM family~\cite{sdofm2024} demonstrate that masked-token
pretraining transfers usefully to physical fields outside the
natural-image domain. The corpus sizes involved
($\mathcal{O}(10^{3})$ to $\mathcal{O}(10^{6})$ samples) define the
regime against which the present work's twenty-one cubes must be
read.

\section{Joint-embedding predictive architectures}
\label{sec:bg-jepa}

The joint-embedding predictive architecture (JEPA) framework,
articulated by LeCun~\cite{lecun2022path}, replaces the
pixel-space reconstruction objective of the masked-autoencoder
family with prediction in the embedding space of a second encoder.
Two encoders are maintained: a context encoder that consumes a
masked view, and a target encoder updated by exponential moving
average of the context encoder's parameters. A small predictor maps
context embeddings to target embeddings; smooth-$L_{1}$ or
$L_{2}$ in the embedding space is the only training signal. No
contrastive negatives are required; representational collapse is
prevented by the asymmetry between the two streams.

The image-JEPA variant (I-JEPA) of
Assran~et~al.~\cite{assran2023ijepa} establishes the paradigm at
scale on ImageNet; the V-JEPA family~\cite{assran2024vjepa,
vjepa2_2025,vjepa21_2026} extends it to video. V-JEPA-2-AC, the
action-conditioned variant of V-JEPA-2, adds a block-causal
predictor and is the architectural template adapted in
Chapter~\ref{ch:architecture}.

The MC-JEPA variant of Bardes~et~al.~\cite{bardes2023mcjepa}
demonstrates that the JEPA encoder is compatible with multiple
downstream heads sharing a single frozen backbone; this property
is used implicitly in Chapter~\ref{ch:probe} where the spatial-mean
winding-flux probe attaches to the frozen encoder without
modifying the pretraining objective. LeJEPA of
Balestriero~et~al.~\cite{balestriero2025lejepa} establishes
theoretical guarantees on JEPA's anti-collapse mechanism.

\section{Solar foundation models}
\label{sec:bg-surya}

Two solar foundation-model lines exist at the time of writing.
Surya, the NASA--IBM model of
Schmude~et~al.~\cite{surya2025,suryabench2025}, pretrains by
forecasting one frame ahead on nine years of twelve-channel SDO
data at $4096^{2}$ spatial resolution. Surya reaches a downstream
flare TSS of $0.436$, the strongest published pixel-only result.
The SDO-FM family~\cite{sdofm2024} explores comparable masked-
autoencoder objectives on the same data source. Both lines
operate at corpus sizes three to four orders of magnitude larger
than the present active-region collection.

The Surya backbone, HelioSpectFormer, is hard-locked to
$4096 \times 4096$ inputs at sixty-minute cadence with thirteen
channels (see Section~\ref{sec:intro-prior}). A LoRA-based
adaptation~\cite{hu2021lora} of Surya to active-region cubes
of variable spatial extent, twelve-minute cadence, and a single
channel was therefore architecturally infeasible. This
incompatibility motivates the present from-scratch path.

\section{Joint-embedding versus adversarial inpainting}
\label{sec:bg-vs-pathak}

The context-encoder family of
Pathak~et~al.~\cite{pathak2016contextencoders}, an adversarial
inpainting approach, was the historical precursor to the masked
autoencoder. Two arguments motivate the choice of
joint-embedding prediction over adversarial inpainting in the
present work.

First, the present field is a noisy pseudoscalar with a
heavy-tailed pixel distribution (\S\ref{sec:data-clip}).
Adversarial training in pixel space rewards visually plausible
reconstructions even at the price of high-frequency fidelity to
the underlying physics; a discriminator trained against such a
field tends to lock onto the heavy-tail outliers rather than the
load-bearing structure. The embedding-space objective discards
high-frequency content by construction, because the predictor
must reproduce only what the target encoder retains.

Second, the V4 pixel-space forecaster trained on the same corpus
exhibited a persistence collapse: the output decayed to the
input frame as training progressed. The pixel-MSE forecaster
chooses to predict the conditional mean of a noisy field, which
is dominated over short horizons by the most recent observation;
the discriminator does not prevent this because the persistence
prediction is locally indistinguishable from a plausible field.
The joint-embedding objective addresses the failure mode at its
root, by predicting in a space the encoder has chosen.

The dense-prediction decoder family
(DPT~\cite{ranftl2021dpt}, SegFormer~\cite{xie2021segformer},
Earthformer~\cite{earthformer2022}) is compatible with a frozen
JEPA encoder and is the natural choice for any subsequent
pixel-space reconstruction head; this is recorded as future work
in Chapter~\ref{ch:summary}.
